\documentclass[11pt]{article}

\usepackage[margin=1in]{geometry}
\usepackage[T1]{fontenc}
\usepackage[utf8]{inputenc}
\usepackage{lmodern}
\usepackage{microtype}
\usepackage{hyperref}
\usepackage{xcolor}
\usepackage{enumitem}
\usepackage{fancyhdr}
\usepackage{longtable}
\usepackage{array}
\usepackage{listings}

\hypersetup{
  colorlinks=true,
  linkcolor=blue,
  urlcolor=blue,
  pdftitle={SpaceAGORA Autonomy and GNC Interface Control Document},
  pdfauthor={SpaceAGORA Interface Inspection}
}

\pagestyle{fancy}
\fancyhf{}
\lhead{SpaceAGORA Autonomy and GNC ICD}
\rhead{\thepage}

\setlist[itemize]{leftmargin=1.4em}
\setlist[enumerate]{leftmargin=1.6em}

\lstdefinestyle{code}{
  basicstyle=\ttfamily\small,
  breaklines=true,
  frame=single,
  columns=fullflexible,
  keepspaces=true,
  showstringspaces=false,
  backgroundcolor=\color{gray!6}
}
\lstset{style=code}

\title{SpaceAGORA Autonomy and GNC\\Interface Control Document}
\author{SpaceAGORA Interface Inspection}
\date{\today}

\begin{document}
\maketitle
\tableofcontents
\newpage

\section{Purpose}

This document captures the code interfaces needed to use SpaceAGORA as the
physics engine for decision-making software stacks, GNC autonomy stacks,
navigation and estimation systems, control systems, and custom dynamics models.

The intended integration pattern is:

\begin{lstlisting}
External autonomy stack
  -> NavigationModel callbacks for sensors and estimators
  -> GuidanceModel callbacks for decisions and mission policy
  -> ControlModel callbacks for command execution and actuator behavior
  -> DynamicsModel effectors only for physical forces and torques
  -> run_simulation(SimulationConfiguration)
\end{lstlisting}

Guidance should not mutate dynamics directly. Guidance/control coupling should
cross typed command boundaries such as \texttt{PropulsiveManeuverCommand} and
\texttt{AerobrakingControlCommand}.

\section{Top-Level Simulation Boundary}

The primary execution entry point is:

\begin{lstlisting}
SpaceAGORA.run_simulation(config::SimulationConfiguration)
\end{lstlisting}

The \texttt{SimulationConfiguration} object is the top-level interface object
that binds the physics environment, spacecraft, GNC callback models, runtime
settings, initial time, and solver tolerances.

Relevant source:
\begin{itemize}
  \item \texttt{src/SpaceAGORA.jl}
  \item \texttt{src/core/state/simulation\_configuration.jl}
\end{itemize}

\texttt{SimulationConfiguration} contains:

\begin{lstlisting}
file_paths
simulation_settings
mission_configuration
environment_model
dynamics_model
guidance_model
navigation_model
control_model
initial_time
integration_tolerances
\end{lstlisting}

\section{Simulation Configuration Interfaces}

\subsection{EnvironmentModel}

\texttt{EnvironmentModel} wires the physical environment:

\begin{lstlisting}
EnvironmentModel(
    planet,
    EI,
    density_model,
    ephemerides_model,
    topography,
    topo_degree,
    topo_order,
    wind,
    thermal_model,
)
\end{lstlisting}

Use this for planet/body selection, atmospheric model selection, ephemerides,
topography, wind enablement, and thermal model selection.

\subsection{DynamicsModel}

\texttt{DynamicsModel} wires spacecraft and physical force/torque effectors:

\begin{lstlisting}
DynamicsModel([spacecraft], (dynamic_effector_1, dynamic_effector_2))
\end{lstlisting}

Use this only for physical dynamics contributors such as gravity, atmospheric
forces, solar radiation pressure, perturbations, and custom physical force or
torque models.

\subsection{NavigationModel}

\texttt{NavigationModel} wires sensor and estimator callbacks:

\begin{lstlisting}
NavigationModel(
    navigation_effectors=(nav_model,),
    navigation_rates=[rate_seconds],
)
\end{lstlisting}

Use this for sensors, estimators, relative navigation, inter-agent measurement
models, communication graph state, and navigation data products.

\subsection{GuidanceModel}

\texttt{GuidanceModel} wires decision and mission-policy callbacks:

\begin{lstlisting}
GuidanceModel(
    guidance_effectors=(guidance_model,),
    guidance_rates=[rate_seconds],
)
\end{lstlisting}

Use this for mission decisions, mode logic, trajectory planning, maneuver
planning, aerobraking policy selection, and autonomy decision outputs.

\subsection{ControlModel}

\texttt{ControlModel} wires control and actuation callbacks:

\begin{lstlisting}
ControlModel(
    control_effectors=(controller,),
    control_rates=[rate_seconds],
)
\end{lstlisting}

Use this for command execution, actuator allocation, reaction wheel control,
thruster scheduling, attitude control, force/torque actuation, and propellant
mass flow.

\section{Spacecraft Object and Motion Configuration}

SpaceAGORA represents each simulated vehicle with a \texttt{SpacecraftModel}.
The spacecraft object defines the vehicle structure, mass properties, hardware
attachment points, and initial translational/rotational state. The spacecraft
does not move by itself; motion is produced by the integrator from the
spacecraft initial condition plus the forces, torques, control outputs, and
mission settings configured around it.

Relevant source:
\begin{itemize}
  \item \texttt{src/vehicle/spacecraft/model.jl}
  \item \texttt{docs/src/user/simulation\_configuration.md}
  \item \texttt{examples/AGORA\_Earth.jl}
  \item \texttt{examples/Earth\_Thruster\_Test.jl}
  \item \texttt{examples/Earth\_RW\_Test.jl}
\end{itemize}

\subsection{SpacecraftModel}

The spacecraft container is:

\begin{lstlisting}
SpacecraftModel(
    joints,
    links,
    root,
    instant_actuation,
    prop_mass,
    inertia_tensor,
    n_reaction_wheels,
    n_thrusters,
    initial_condition,
    id,
)
\end{lstlisting}

The key fields are:

\begin{longtable}{>{\ttfamily}p{0.26\linewidth}p{0.66\linewidth}}
\textbf{Field} & \textbf{Role} \\
\hline
joints & Structural joints connecting links in a multi-body vehicle model. \\
links & List of body links. The constructor ensures the root link is included. \\
root & Main bus or core body link. \\
instant\_actuation & Whether control inputs such as panel angle changes are applied instantly. \\
dry\_mass & Computed by summing link masses in the constructor. \\
prop\_mass & Propellant mass available for maneuvers. \\
inertia\_tensor & Spacecraft body-frame inertia tensor used by attitude dynamics. \\
n\_reaction\_wheels & Number of reaction wheels represented in the model. \\
n\_thrusters & Number of thrusters represented in the model. \\
initial\_condition & Orbit, Cartesian state, attitude, and angular-rate initial condition. \\
id & Spacecraft identifier for multi-spacecraft simulations. \\
\end{longtable}

A minimal single-body spacecraft can be constructed as:

\begin{lstlisting}
root = Link(
    root=true,
    dims=[1.0, 1.0, 1.0],
    ref_area=1.0,
    m=100.0,
)

ic = InitialCondition(
    ra=planet.Rp_e + 700e3,
    rp=planet.Rp_e + 400e3,
    i=28.5,
    omega=10.0,
    Omega=20.0,
    nu=0.0,
)

spacecraft = SpacecraftModel(
    root=root,
    links=[root],
    prop_mass=20.0,
    initial_condition=ic,
    id=1,
)
\end{lstlisting}

In the actual Julia source, the angular orbital-element names use Greek
identifiers \texttt{omega}, \texttt{Omega}, and \texttt{nu} in examples when
rendered here for portability; repository examples may use the corresponding
Unicode symbols directly.

\subsection{Link}

\texttt{Link} represents a body element such as a bus, panel, appendage, or
component body. Important fields include:

\begin{longtable}{>{\ttfamily}p{0.26\linewidth}p{0.66\linewidth}}
\textbf{Field} & \textbf{Role} \\
\hline
root & Marks the main bus/core body. \\
r, q & Link position and attitude. For the root, these are inertial/root states; for non-root links they are body-relative. \\
rdot, omega & Link velocity and angular velocity. \\
dims & Box dimensions used for simple geometry and inertia defaults. \\
ref\_area & Reference area used by aerodynamic models. \\
m & Link mass. \\
mass & Link mass matrix. \\
inertia & Link inertia matrix. \\
a\_b, b\_b & Body-frame left/right extents used for geometry and joints. \\
alpha, beta, theta & Aerodynamic attitude/flow angles. \\
reflection\_coefficient & Surface optical/aerodynamic reflection coefficient. \\
rw\_assembly & Reaction wheel assembly state and limits. \\
net\_force, net\_torque & Per-link accumulated force and torque buffers. \\
attitude\_control\_rate & Link-level attitude control cadence where used. \\
SRP\_facets & Facets for solar radiation pressure models. \\
J\_thruster & Thruster allocation/Jacobian matrix. \\
thrusters & Thruster hardware attached to the link. \\
magnets & Magnetic dipoles attached to the link. \\
\end{longtable}

The constructor computes default mass and inertia values from mass and box
dimensions when explicit matrices are not supplied.

\subsection{Joint}

\texttt{Joint} connects two \texttt{Link} objects for multi-body spacecraft.
It stores body-frame connection points, translational/rotational stiffness,
damping, and displacement state:

\begin{lstlisting}
Joint(
    link1,
    link2;
    p1=link1.b_b,
    p2=link2.a_b,
    Kx=...,
    Kt=...,
    Cx=...,
    Ct=...,
)
\end{lstlisting}

Use joints when the spacecraft model needs articulated panels, flexible
appendages, or multi-body geometry that is more detailed than a single rigid
root link.

\subsection{Initial Conditions}

SpaceAGORA supports Keplerian and Cartesian initial conditions.

Keplerian mode:

\begin{lstlisting}
ic = InitialCondition(
    ra=planet.Rp_e + 1200e3,   # apoapsis radius from planet center, m
    rp=planet.Rp_e + 400e3,    # periapsis radius from planet center, m
    i=28.5,                    # inclination, deg
    omega=10.0,                # argument of periapsis, deg
    Omega=20.0,                # right ascension of ascending node, deg
    nu=0.0,                    # true anomaly, deg
    q=q0,                      # initial attitude quaternion
    ang_vel=w0,                # initial body angular velocity, rad/s
)
\end{lstlisting}

The keyword constructor also supports semi-major axis and eccentricity:

\begin{lstlisting}
ic = InitialCondition(
    a=planet.Rp_e + 800e3,
    e=0.001,
    i=28.5,
    omega=10.0,
    Omega=20.0,
    nu=0.0,
)
\end{lstlisting}

Cartesian mode:

\begin{lstlisting}
ic = CartesianInitialCondition(
    pos=[6_778_137.0, 0.0, 0.0],  # inertial position, m
    vel=[0.0, 7784.0, 0.0],       # inertial velocity, m/s
    q=q0,
    ang_vel=w0,
)
\end{lstlisting}

All Keplerian angular inputs are accepted in degrees and converted internally
to radians.

\subsection{How Spacecraft Motion Is Produced}

The propagation path is:

\begin{enumerate}
  \item The spacecraft initial condition seeds the integrator state.
  \item \texttt{MissionConfiguration} selects mission duration, orbit count,
        Keplerian/two-phase mode, output cadence, and whether attitude dynamics
        are enabled.
  \item \texttt{EnvironmentModel} provides planet, atmosphere, ephemerides,
        topography, wind, and thermal context.
  \item \texttt{DynamicsModel.dynamic\_effectors} contributes physical forces
        and torques through \texttt{wrench} or \texttt{calcForceTorque}.
  \item \texttt{ControlModel.control\_effectors} contributes commanded forces,
        torques, and mass flow through \texttt{calcControlForceTorque} and
        \texttt{calcControlMassFlowRate}.
  \item The ODE RHS sums physical and control forces/torques and advances the
        spacecraft translational state, mass state, heat state, and, when
        enabled, attitude state.
\end{enumerate}

Translational motion is controlled by initial position/velocity and the net
force model. Rotational motion is controlled by initial attitude/angular rate,
net body torque, inertia, and the \texttt{orientation\_sim} mission setting.

\subsection{Translational Motion Configuration}

To configure where the spacecraft goes, configure:

\begin{lstlisting}
spacecraft.initial_condition
SimulationConfiguration.environment_model
SimulationConfiguration.dynamics_model
SimulationConfiguration.control_model
MissionConfiguration(
    mission_type=MissionTime,
    mission_time=...,
    keplerian=true_or_false,
    orientation_sim=false,
)
\end{lstlisting}

Example:

\begin{lstlisting}
dynamics = DynamicsModel(
    [spacecraft],
    (
        InverseSquaredJ2GravityModel(),
        AerodynamicCoefficientConstant(...),
    ),
)
\end{lstlisting}

The physical effectors determine truth acceleration. Navigation does not move
the truth state.

\subsection{Rotational Motion Configuration}

To propagate attitude, set:

\begin{lstlisting}
MissionConfiguration(
    orientation_sim=true,
    ...
)
\end{lstlisting}

Then provide an initial quaternion and angular rate in the spacecraft initial
condition:

\begin{lstlisting}
q0 = normalize(SVector{4, Float64}(0.2, -0.1, 0.15, 0.95))
w0 = SVector{3, Float64}(0.003, -0.004, 0.005)

ic = InitialCondition(a, e, i, omega, Omega, nu, q0, w0)
\end{lstlisting}

Rotational motion responds to body-frame torques from dynamics and control
effectors, as well as spacecraft inertia and link-level geometry where the
selected models use it.

\subsection{Navigation, Guidance, and Control Effects on Motion}

Navigation callbacks observe and estimate state. They should update estimator,
sensor, or communication-model state, but they should not directly advance the
truth state.

Guidance callbacks make decisions and produce commands or plans. They should
not directly mutate physical dynamics. For aerobraking and propulsive maneuver
logic, use typed command boundaries such as:

\begin{lstlisting}
PropulsiveManeuverCommand
PropulsiveBurnPlan
AerobrakingControlCommand
\end{lstlisting}

Control callbacks execute commands. Motion changes when control effectors
return nonzero force, torque, or mass flow:

\begin{lstlisting}
calcControlForceTorque(model, u, p, sat_idx, t)
calcControlMassFlowRate(model, u, p, sat_idx, t)
\end{lstlisting}

\subsection{Minimal End-to-End Spacecraft Setup}

The following is the minimal shape of a spacecraft simulation setup:

\begin{lstlisting}
using SpaceAGORA
const SM = SpaceAGORA.SimulationModel

planet = SM.make_no_gram_planet(:earth)

root = SM.Link(
    root=true,
    dims=[1.0, 1.0, 1.0],
    ref_area=1.0,
    m=100.0,
)

ic = SM.InitialCondition(
    ra=planet.Rp_e + 700e3,
    rp=planet.Rp_e + 400e3,
    i=28.5,
    omega=10.0,
    Omega=20.0,
    nu=0.0,
)

spacecraft = SM.SpacecraftModel(
    root=root,
    links=[root],
    prop_mass=20.0,
    initial_condition=ic,
    id=1,
)

config = SM.SimulationConfiguration(
    mission_configuration=SM.MissionConfiguration(
        mission_type=SM.MissionTime,
        mission_time=3600.0,
        keplerian=true,
        orientation_sim=false,
        number_of_orbits=1,
        num_steps_to_save=1000,
        data_rate=10.0,
    ),
    environment_model=SM.EnvironmentModel(
        planet=planet,
        EI=120.0,
        density_model=SM.NoAtmosphereModel(),
        ephemerides_model=SM.SimpleEphemeridesModel(),
        thermal_model=SM.MaxwellianHeat(
            thermal_accomodation_factor=1.0,
            planet=planet,
        ),
        topography=false,
        wind=false,
    ),
    dynamics_model=SM.DynamicsModel(
        [spacecraft],
        (SM.InverseSquaredJ2GravityModel(),),
    ),
    guidance_model=SM.GuidanceModel(
        guidance_effectors=(),
        guidance_rates=Float64[],
    ),
    navigation_model=SM.NavigationModel(
        navigation_effectors=(),
        navigation_rates=Float64[],
    ),
    control_model=SM.ControlModel(
        control_effectors=(),
        control_rates=Float64[],
    ),
    initial_time=SM.InitialTime(
        year=2024,
        month=1,
        day=1,
        hour=0,
        minute=0,
        second=0.0,
    ),
    integration_tolerances=SM.IntegrationTolerances(),
)

sol = SpaceAGORA.run_simulation(config; return_solution=true)
\end{lstlisting}

\section{Dynamic Physics Effector Interface}

For new physics models, prefer the typed \texttt{wrench} interface.

\begin{lstlisting}
struct MyEffector <: SpaceAGORA.AbstractForceTorqueModel
    # model parameters
end

SpaceAGORA.environment_requirements(model::MyEffector) =
    SpaceAGORA.EffectorEnvironmentRequirements(
        planet_frame=true,
        atmosphere=false,
        solar=false,
        third_body_names=(),
    )

function SpaceAGORA.wrench(
    model::MyEffector,
    x::SpaceAGORA.StateSample,
    env::SpaceAGORA.EnvironmentSample,
    t::Float64,
)
    force_ii = ...
    torque_body = ...
    return force_ii, torque_body
end
\end{lstlisting}

The returned force is inertial-frame SI force. The returned torque is body-frame
SI torque.

Relevant source:
\begin{itemize}
  \item \texttt{src/core/types/effector\_sampling.jl}
  \item \texttt{src/simulation/engine/dynamics\_rhs.jl}
  \item \texttt{templates/force\_torque\_model\_template.jl}
\end{itemize}

\subsection{StateSample}

\texttt{StateSample} provides:

\begin{lstlisting}
pos_ii       # inertial position, m
vel_ii       # inertial velocity, m/s
mass_kg      # spacecraft mass, kg
q_ib         # inertial-to-body attitude quaternion, when available
omega_body   # body-frame angular rate, when available
spacecraft   # typed spacecraft model handle
\end{lstlisting}

\subsection{EnvironmentSample}

\texttt{EnvironmentSample} provides optional sampled environment fields
requested by \texttt{environment\_requirements}:

\begin{lstlisting}
planet
planet_frame
atmosphere
solar
third_bodies
\end{lstlisting}

\subsection{Environment Requirements}

Use \texttt{EffectorEnvironmentRequirements} to request only the sampled data
required by a physics effector:

\begin{lstlisting}
EffectorEnvironmentRequirements(
    planet_frame=true,
    atmosphere=true,
    solar=false,
    third_body_names=("Sun", "Moon"),
)
\end{lstlisting}

\subsection{Solver Partitioning}

Optional solver partition hooks:

\begin{lstlisting}
SpaceAGORA.solver_partition(model) -> :explicit | :implicit
SpaceAGORA.gravity_backbone_structure(model) ->
    :unsupported | :position_only_static_gravity
SpaceAGORA.gravity_backbone_acceleration_ii(model, x, env, t)
SpaceAGORA.gravity_backbone_kick_structure(model) ->
    :unsupported | :velocity_kick_explicit
SpaceAGORA.gravity_backbone_kick_acceleration_ii(model, x, env, t)
\end{lstlisting}

Use these only when a dynamics model must participate in specialized solver
paths.

\subsection{Legacy Force/Torque Hook}

The legacy hook still exists:

\begin{lstlisting}
SpaceAGORA.calcForceTorque(model, x, p, sat_idx) ->
    (force_n, torque_n_m)
\end{lstlisting}

Use \texttt{wrench} for new integrations where possible.

\section{Atmosphere and Density Model Interface}

Custom density models subtype:

\begin{lstlisting}
SpaceAGORA.AbstractDensityModel
\end{lstlisting}

Required scalar hook:

\begin{lstlisting}
SpaceAGORA.getDensity(model, h, lat, lon, el_time, wind[, p]) ->
    (rho, temperature, wind_vec)
\end{lstlisting}

Optional batch hook:

\begin{lstlisting}
SpaceAGORA.getDensityBatch!(
    rhos,
    Ts,
    winds,
    model,
    hs,
    lats,
    lons,
    el_time,
    wind,
    p,
)
\end{lstlisting}

Register the model as:

\begin{lstlisting}
environment_model.density_model
\end{lstlisting}

Relevant source:
\begin{itemize}
  \item \texttt{src/SpaceAGORA.jl}
  \item \texttt{templates/density\_model\_template.jl}
\end{itemize}

\section{Navigation and Estimation Interface}

Navigation models are periodic callback models. They are invoked once per
configured navigation rate for each spacecraft index.

Implement:

\begin{lstlisting}
SpaceAGORA.SimulationModel.calcNavigationEffect!(
    model,
    u,
    p,
    t::Float64,
    sat_idx::Int,
)
\end{lstlisting}

Register:

\begin{lstlisting}
NavigationModel(
    navigation_effectors=(model,),
    navigation_rates=[rate_seconds],
)
\end{lstlisting}

The callback receives:

\begin{lstlisting}
model    # user navigation model instance
u        # integrator state
p        # ODE/runtime parameters
t        # simulation time, s
sat_idx  # spacecraft index
\end{lstlisting}

Appropriate responsibilities:
\begin{itemize}
  \item sensor simulation
  \item measurement generation
  \item estimator updates
  \item relative navigation
  \item covariance tracking
  \item inter-agent observation state
  \item communication graph updates
\end{itemize}

Relevant source:
\begin{itemize}
  \item \texttt{src/gnc/navigation/navigation\_hooks.jl}
  \item \texttt{src/simulation/callbacks/navigation\_guidance\_callbacks.jl}
  \item \texttt{examples/Earth\_Navigation.jl}
\end{itemize}

Note: this hook exists through \texttt{SpaceAGORA.SimulationModel}; it is not
promoted to the root API in the same way as \texttt{wrench} and control hooks.

\section{Guidance and Decision-Making Interface}

Guidance models are periodic callback models. They are invoked once per
configured guidance rate for each spacecraft index.

Implement:

\begin{lstlisting}
SpaceAGORA.SimulationModel.calcGuidanceEffect!(
    model,
    u,
    p,
    t::Float64,
    sat_idx::Int,
)
\end{lstlisting}

Register:

\begin{lstlisting}
GuidanceModel(
    guidance_effectors=(model,),
    guidance_rates=[rate_seconds],
)
\end{lstlisting}

Appropriate responsibilities:
\begin{itemize}
  \item mission policy decisions
  \item mode selection
  \item maneuver planning
  \item trajectory guidance
  \item target selection
  \item aerobraking strategy selection
  \item command generation for downstream control
\end{itemize}

Relevant source:
\begin{itemize}
  \item \texttt{src/gnc/guidance/guidance\_hooks.jl}
  \item \texttt{src/simulation/callbacks/navigation\_guidance\_callbacks.jl}
\end{itemize}

Note: this hook exists through \texttt{SpaceAGORA.SimulationModel}; it is not
promoted to the root API in the same way as \texttt{wrench} and control hooks.

\section{Control and Actuation Interface}

Control models use the stable root API:

\begin{lstlisting}
struct MyController <: SpaceAGORA.AbstractControlEffectorModel
    # controller state and parameters
end

SpaceAGORA.calcControlEffect!(model, u, p, t::Float64, sat_idx::Int)
SpaceAGORA.calcControlForceTorque(model, u, p, sat_idx::Int, t::Float64)
SpaceAGORA.calcControlMassFlowRate(model, u, p, sat_idx::Int, t::Float64)
\end{lstlisting}

Register:

\begin{lstlisting}
ControlModel(
    control_effectors=(model,),
    control_rates=[rate_seconds],
)
\end{lstlisting}

The control callback path periodically calls \texttt{calcControlEffect!}. The
dynamics RHS accumulates \texttt{calcControlForceTorque} and
\texttt{calcControlMassFlowRate}.

Appropriate responsibilities:
\begin{itemize}
  \item actuator command execution
  \item attitude control
  \item reaction wheel control
  \item thruster control
  \item burn scheduling
  \item force/torque generation
  \item propellant mass depletion
\end{itemize}

Relevant source:
\begin{itemize}
  \item \texttt{src/SpaceAGORA.jl}
  \item \texttt{src/simulation/callbacks/control\_callbacks.jl}
  \item \texttt{src/simulation/engine/dynamics\_rhs.jl}
  \item \texttt{templates/control\_hook\_template.jl}
  \item \texttt{examples/Earth\_Thruster\_Test.jl}
\end{itemize}

\section{Typed Command Interfaces}

SpaceAGORA defines typed command structs for GNC/control boundaries:

\begin{lstlisting}
PropulsiveManeuverCommand
PropulsiveBurnPlan
AerobrakingControlCommand
\end{lstlisting}

Relevant source:
\begin{itemize}
  \item \texttt{src/gnc/command\_types.jl}
\end{itemize}

\subsection{PropulsiveManeuverCommand}

Fields:

\begin{lstlisting}
valid::Bool
delta_v_mps::Float64
direction_rad::Float64
source_orbit::Int64
\end{lstlisting}

\subsection{PropulsiveBurnPlan}

Fields:

\begin{lstlisting}
valid::Bool
delta_v_mps::Float64
direction_rad::Float64
source_orbit::Int64
start_burn_s::Float64
stop_burn_s::Float64
thrust_n::Float64
isp_s::Float64
commanded_impulse_n_s::Float64
propellant_required_kg::Float64
\end{lstlisting}

\subsection{AerobrakingControlCommand}

Fields:

\begin{lstlisting}
alpha_command::Float64
\end{lstlisting}

\section{Aerobraking Strategy and Policy Interface}

Aerobraking has a typed strategy/policy boundary.

Canonical interfaces:

\begin{lstlisting}
AerobrakingGuidanceInput
AerobrakingGuidanceOutput
AbstractAerobrakingStrategy
AerobrakingPolicyConfig
AbstractAerobrakingPolicySelector
AerobrakingStrategyKind
E_EDG
T_EDG
\end{lstlisting}

Relevant source:
\begin{itemize}
  \item \texttt{src/gnc/guidance/aerobraking/interfaces.jl}
  \item \texttt{src/mission/operations/aerobraking\_policy/policy\_types.jl}
  \item \texttt{docs/architecture/gnc\_aerobraking\_boundary\_contract.md}
\end{itemize}

Policy selection:

\begin{lstlisting}
select_strategy(selector, config, input) -> AerobrakingStrategyKind
\end{lstlisting}

Guidance dispatch:

\begin{lstlisting}
dispatch_aerobraking_guidance(kind, input)
dispatch_aerobraking_guidance(selector, config, input)
\end{lstlisting}

The repository includes:

\begin{lstlisting}
DefaultAerobrakingPolicySelector
DRLPolicyAdapterStub
\end{lstlisting}

\texttt{DRLPolicyAdapterStub} intentionally throws \texttt{Not implemented}. A
reinforcement learning or decision-making adapter should implement the
\texttt{AbstractAerobrakingPolicySelector} contract rather than reaching into
control or dynamics internals.

\section{Runtime Scheduling Semantics}

Navigation, guidance, and control callbacks are scheduled from their model
containers and rates:

\begin{lstlisting}
NavigationModel.navigation_effectors
NavigationModel.navigation_rates

GuidanceModel.guidance_effectors
GuidanceModel.guidance_rates

ControlModel.control_effectors
ControlModel.control_rates
\end{lstlisting}

Each callback is invoked over all spacecraft indices:

\begin{lstlisting}
for sat_idx in 1:num_sats
    callback(model, integrator.u, integrator.p, integrator.t, sat_idx)
end
\end{lstlisting}

Control effectors that are \texttt{BaseThrusterModel} receive event-driven burn
scheduling behavior and tstop registration.

\section{Public API Stability Notes}

The strongest stable root interfaces are:

\begin{lstlisting}
SpaceAGORA.run_simulation
SpaceAGORA.AbstractForceTorqueModel
SpaceAGORA.wrench
SpaceAGORA.environment_requirements
SpaceAGORA.AbstractDensityModel
SpaceAGORA.getDensity
SpaceAGORA.AbstractControlEffectorModel
SpaceAGORA.calcControlEffect!
SpaceAGORA.calcControlForceTorque
SpaceAGORA.calcControlMassFlowRate
\end{lstlisting}

Navigation and guidance callbacks are present and used by examples and tests,
but currently live under:

\begin{lstlisting}
SpaceAGORA.SimulationModel.calcNavigationEffect!
SpaceAGORA.SimulationModel.calcGuidanceEffect!
\end{lstlisting}

For an external autonomy package, isolate those calls behind a small adapter
layer so future internal reshuffles in SpaceAGORA are localized.

\section{Recommended External Stack Shape}

Use this layering:

\begin{lstlisting}
Autonomy application
  - Owns high-level mission behavior
  - Owns policy/ML/RL decision models
  - Owns autonomy state machines

SpaceAGORA adapter layer
  - Builds SimulationConfiguration
  - Wraps navigation/guidance/control hooks
  - Converts autonomy commands into typed SpaceAGORA command structs
  - Converts SpaceAGORA state into autonomy observations

SpaceAGORA
  - Owns propagation
  - Owns force/torque accumulation
  - Owns environment sampling
  - Owns callback scheduling
  - Owns output and solver execution
\end{lstlisting}

Use \texttt{DynamicsModel} for physical truth. Use
\texttt{NavigationModel}, \texttt{GuidanceModel}, and \texttt{ControlModel}
for autonomy software behavior.

\end{document}
